% auto-generated by make exhibits -- do not edit
\begin{table}[!htbp]\centering\footnotesize
\caption{Demography of blocked majorities (descriptive; California)}\label{tab:blocked}
\begin{tabular}{@{}>{\raggedright\arraybackslash}p{0.237\linewidth}>{\centering\arraybackslash}p{0.115\linewidth}>{\centering\arraybackslash}p{0.123\linewidth}>{\centering\arraybackslash}p{0.115\linewidth}>{\centering\arraybackslash}p{0.164\linewidth}>{\centering\arraybackslash}p{0.147\linewidth}@{}}
\toprule
covariate & blocked & cleared & sub-50 & blocked-cleared (SE) & blocked-sub50 (SE) \\
\midrule
renter share & 0.370 (n=277) & 0.408 (n=1526) & 0.351 (n=161) & -0.038 (0.006) & +0.019 (0.010) \\
Gini (B19083) & 0.442 (n=277) & 0.450 (n=1526) & 0.437 (n=161) & -0.008 (0.002) & +0.005 (0.004) \\
fractionalization & 0.538 (n=277) & 0.558 (n=1526) & 0.511 (n=161) & -0.020 (0.010) & +0.027 (0.014) \\
median HH income \$ & 64,710 (n=277) & 68,508 (n=1526) & 68,344 (n=161) & -3,798 (1,463) & -3,635 (2,284) \\
65+ share & 0.130 (n=277) & 0.130 (n=1526) & 0.142 (n=161) & +0.000 (0.003) & -0.012 (0.005) \\
SAIPE child poverty (schools) & 0.155 (n=107) & 0.175 (n=696) & 0.142 (n=95) & -0.020 (0.010) & +0.013 (0.012) \\
\bottomrule
\end{tabular}
\begin{minipage}{0.96\linewidth}\vspace{2pt}\scriptsize
\emph{Notes:} Blocked = majority short of the supermajority threshold; within-matched ACS/SAIPE at the D5 grain ladder; DESCRIPTIVE, no causal claim.
\end{minipage}
% BEGIN READING (strip for submission)
\begin{minipage}{0.96\linewidth}\vspace{4pt}\scriptsize
\textbf{Reading.} Blocked majorities arise in less affluent, mid-composition school districts: the incidence of the higher bar.
\end{minipage}
% END READING
\end{table}
